\section{Variación del corrector y coste del cambio de referencia}
Para una matriz $B=P-J$ escribimos $G_L=BB^\top,G_R=B^\top B$,
$S_v=(I-G_v)^{-1}$, $T_v=S_v-I$, $V=S_LB=BS_R$, $U=V-B$.
El vector $j$ reúne las diagonales de $T_v-G_v$, $j_0=\tr(T_R-G_R)$.
El bloque diagonal de $H$ es $(-T_L,-T_R)$ y su bloque superior cruzado es $V$.
Denotemos $\sigma_3(D)=\sum D_{ij}^3$. El funcional de ronda 8 es exactamente
\begin{align*}
 \Xi={}&(j_0^2/2+\tr T_R^2-3j_0/2)/k
 -3j^\top j/4+j^\top Hj/2\\
 &-(\sigma_3(T_L)+\sigma_3(T_R))/3+2\sum_{ij}B_{ij}U_{ij}^2+2\sigma_3(U)/3 .
\end{align*}
\paragraph{Estimación diferencial revisada.}
Para toda entrada independiente $B_{ab}$ del segmento del apartado anterior,
$|\partial_{ab}\Xi|\le C/k^2$. En efecto $B,U,V,T_s,H$ tienen entradas
$O(k^{-1})$ y sumas absolutas por fila y columna acotadas. Para $T_L$ se usa
$B S_RB^\top$ y para $V=B+BT_R$; la brecha y las normas de filas y columnas
$O(k^{-1/2})$ dan las cotas. Para $E=e_ae_b^\top$ se tienen las diferenciales
\[
 dT_L=xv^\top+vx^\top,\quad x=S_Le_a,\quad v=Ve_b,\qquad
 dV=S_LE S_R+VE^\top V.
\]
A la derecha se trasponen los papeles. Cada diferencial de $T,G,U,H$ tiene
suma absoluta de entradas $O(1)$; la diagonal de $dT$ y de $dG$ tiene suma absoluta
$O(k^{-1})$, ya que $\|x\|_1=O(1)$ y $\|v\|_\infty=O(k^{-1})$.
Resultan $|dj_0|+\|dj\|_1=O(k^{-1})$ y $|d\tr T_R^2|=O(k^{-1})$.
En la forma cuadrática, $\|j\|_\infty+\|Hj\|_\infty=O(k^{-1})$.
Para los cubos y $BU^2$ se usan dos factores de tamaño $O(k^{-1})$ y la
suma absoluta de la diferencial del tercero. Esto da la cota para cada término.
La integral sobre el segmento cuesta $C/k^2$ por su cambio total en norma de
suma de entradas. En la matriz con ceros añadidos las sumas de la fórmula son
las del hijo; el paso de $1/k$ a $1/(k-1)$ en la traza cuesta $O(k^{-2})$.
También el término universal cambia en $O(k^{-2})$. Así
\[
 |\Delta u_T|\le C/k^2,\qquad
 |\log D_\tau|\le\sum_{k>K}2C/k^2\le C'/K.
\]
El seguimiento disponible bajo $\nu$ se transfiere en historias detenidas por un
factor entre $e^{-C'/K}$ y $e^{C'/K}$; no se obtiene por ello una probabilidad
de salida mejor que la original.
\paragraph{El predictor con permanente exacto.}
Para $L^\sharp=L-L_T+\log\per A$, sea $\Lambda=L-L_T$, $A_c$ la máscara verdadera
del hijo y $A^-_c$ el menor de torres al borrar sólo fila y columna. Con la ley
$p^{\rm per}_c=\per A^-_c/\per A$ en casillas de la fila se tiene exactamente
\[
 \rho_i^\sharp=\E_{\rm per}
 \left[\frac{\per A_c}{\per A^-_c}\exp\{\Lambda(c)-\Lambda(B)\}\right].
\]
Se suman sólo hijos de masa de torres positiva. Un cociente cero aporta cero
si los factores se definen por los pesos correspondientes. La igualdad sale al
sustituir $L^\sharp$ en la suma de fila. En una máscara completa $3$ por $3$ con
una restricción doble sobre $(1,1),(2,2)$, elegir $(1,1)$ deja dos permutaciones
de torres y sólo una respeta el veto. Es un ejemplo del cociente $1/2$
independiente de la existencia de la rama de Bethe en ese tamaño.
El error terminal para $L^\sharp$ es exactamente $Y_D=g-g_T$.
